

\subsubsubsection{Block-DAG Interaction With The PoR Graph}

\label{sec:dags-and-the-por-graph}

Up to now we have mostly ignored that \UT{} block-DAG simplex-chains exist within a broader context.
Particularly, we should ask: \emph{do the Axioms of \AxOfAvailability and \AxOfMaximalReflection change the game?}

Recall that, when defining the \AxiomOfMaximalReflection, we make sure to recursively include blocks' ancestors, not just reflections.
This was for two reasons.
First, \BL should not explicitly reflect its own ancestors -- links to parents \emph{are} reflections, but are the trivial case where the local and reflecting chains are the same, and the local and reflecting blocks are the same.
Second, \BL should not skip accounting for \BR's implied reflection of its parent(s), so we treat parents like any other reflected block.
We must be sure not to leave any gaps.

There is still one gap, though: \cL blocks that are not an ancestor of \BL.

It should be clear that parent links are qualitatively different from reflections: child blocks attest to the validity of their parent blocks, and to continuity between these two blocks.
% This was the reason for the predicate's (ii)(b) condition
``PoR'' between these blocks is also free, unlike explicit PoRs between blocks of different chains.

Unlike traditional chains, \UT{} blocks are not only allowed to have multiple parents, but \emph{can always include a valid block as an ancestor}.
Recall the \AxiomOfMaximalReflection's principle: \emph{\PrincipleOfMaximalReflection} % space but no full stop
Considering the symmetry between parents and reflected blocks, we can devise a new principle: \emph{A miner should build on all possible chain-tips.}

In effect, we want to ensure that a dishonest miner cannot evade the contradiction in refusing to build on a valid chain-tip whilst acknowledging it as available.
If it is unavailable then the miner should not reflect the blocks that imply it is available; and if it is available and valid then it should be a parent of their draft block.

So far, this principle is fine for honest miners building on \emph{valid} blocks -- including every chain-tip as a parent is what we expect honest miners to do by default.
However, what are honest miners to do if there is an invalid block?

We already know enough to answer this, of course: link back to, or \emph{flag}, that block as invalid with an attached fraud proof, and never use that block as a parent.
The transactions in that block aren't important, so we don't care about keeping them --- the offending transaction is part of the fraud proof in any case.
And we can reconstruct the PoRs sub-block using $O(1)$ data thanks to \AxOfMaximalReflection.
Thus, an invalid block is only as large as its header, list of reflected PoR graph-tips, and fraud proof.
That's $O(\log c)$ in the worst case.

% It is only in the case of an invalid block that it cannot be an ancestor, and in that case the miner \emph{must} link to it with the fraud proof attached.
% This link is sort of like a reflection, but it comes with a note saying \emph{this block will never be valid, and here's why}.
% Since the transactions within such a block are non-canonical, we don't count it as an ancestor.
% Thus, the entire PoR graph learns that the block is invalid.

\subsubsubsection{\AxiomOfRawUnifiedAncestry}

\DeclareAxiom{UnifiedAncestry}{
    A miner should use all valid chain-tips as parents, and flag invalid ones.
}{
    A block, \BL{}, is invalid unless, for all reflected blocks \BR{}:
    \begin{axiomreqs}
        \item \BR{} is valid under this axiom; and
        \item each \cL block reflected by \BR{}:
        \begin{enumerate}[label=\alph*.]
            \item is an ancestor of \BL{}; or
            % \item shares no ancestors with \BL; or
            \item is flagged as invalid by \BL{}; or
            \item was flagged as invalid by an ancestor of \BL{}.
        \end{enumerate}
    \end{axiomreqs}
}

This new axiom changes the game for block-DAGs.
Previously, an attacker could delay transactions for $\sim {A_\text{blocks}}^2$ many blocks because they could evade the contradiction in the PoR graph.
Now, they cannot.
Thus, empty block DoS attacks against simplex-chain block-DAGs are limited to $\sim A_\text{blocks}$ in length (on average).\footnote{
    We will reduce this further in \autoref{sec:miner-resonance}.
}
In terms of DoS potential, this is \emph{exactly} the same as if the simplex-chain were a traditional chain and the attacker always built on the best block, but only made empty blocks.
Now that we've added the \AxiomOfUnifiedAncestry, the alternative for the attacker is producing only invalid blocks.

\aside{
    Note that this axiom applies in the context of the PoR graph rather than local chain-state.\\
    The entire PoR graph learns of and verifies the fraud proof.
}

\subsubsubsection{NIPoPPoWs + NIPoPoWRs = NIPoPPoWaPoWRs}

In \autoref{sec:blockdag-nipopows} we devised NIPoPPoWs, modified NIPoPoWs that work with block-DAGs by accounting for posterity.
That construction depended a \emph{best-blocks} proof in combination with an \emph{on-main only} proof.
In the case of simplex-chain block-DAGs, we do not want to use the same best-blocks proof as it only accounts for local chain-work and does not include reflections.

Instead of a standard best-blocks NIPoPoWs, we can use NIPoPoWRs from \autoref{sec:nipopowrs}.
Functionally, both serve the same purpose: proving that a block \emph{exists}.

However, the best-blocks proof had substantial overlap with the on-main only proof -- that isn't the case for NIPoPoWRs as \emph{most} of the \mu-superblocks will be from reflecting chains.
Rater, we expect the best posteritorial \mu-superchain to have a \mu level approximately $\log_2({N_1})$ lower than that of the best PoR \mu-superDAG.



\subsubsubsubsection{Integrating Fraud Proofs}
